\chapter{Planteo del Problema y Selección de Topología}
\label{chap:seleccion_topologia}

En este capítulo se analiza en detalle la situación práctica planteada por la cátedra, justificando desde el
punto de vista de la ingeniería electrónica la topología óptima para satisfacer las especificaciones del sistema.

\section{Descripción del Requerimiento Práctico}

El sistema a desarrollar tiene como propósito acondicionar la señal generada por un micrófono piezoeléctrico
ubicado a una distancia de $30\unit{\meter}$ del equipamiento de grabación de audio. Los parámetros de diseño
suministrados por la consigna son:
\begin{itemize}
  \item Nivel de señal del transductor: $v_s = 200\mV_{\mathrm{rms}}$.
  \item Resistencia interna del micrófono: $R_S \approx 10\kohm$.
  \item Amplitud máxima admisible por el equipo de grabación: $v_{L,\mathrm{max}} = 2\Vpp$.
  \item Impedancia de entrada del equipo de grabación (carga): $R_L \approx 10\kohm$.
  \item Rango de frecuencias de operación: $100\Hz$ a $20\kHz$ (espectro audible).
  \item Longitud de línea de transmisión: $30\unit{\meter}$ de cable apantallado.
\end{itemize}

\section{Análisis de Amplitud y Ancho de Banda}

La tensión senoidal provista por el micrófono presenta un valor pico a pico de:
\begin{equation}
  v_{s,\mathrm{pp}} = 2\sqrt{2} \cdot v_s = 2\sqrt{2} \cdot 200\mV \approx 565.69\mVpp \approx 0.57\Vpp
  \label{eq:vs_pp}
\end{equation}

Dado que el equipo de grabación admite una señal máxima de $2\Vpp$ y la señal del micrófono es de $0.57\Vpp$,
se concluye que la etapa no debe aportar ganancia de tensión significativa ($A_v \approx 1$). Si se introdujera
una ganancia de tensión moderada (por ejemplo $A_v \ge 4$), la amplitud superaría los $2\Vpp$ tolerados por el
sistema receptor.

Por otra parte, la conexión directa del micrófono al equipo de grabación sin circuito de adaptación presentaría
dos problemas severos:
\begin{enumerate}
  \item Atenuación resistiva por divisor de tensión: debido a que $R_S = R_L = 10\kohm$, la mitad de la tensión
    se perdería internamente en el transductor, reduciendo la amplitud a $0.28\Vpp$.
  \item Filtro pasabajos por capacidad del cable: una línea coaxial de audio típica presenta una capacidad distribuida
    del orden de $100\unit{\pico\farad\per\meter}$, lo que representa una capacidad parásita total de
    $C_{\mathrm{cable}} \approx 30\unit{\meter} \times 100\pF = 3\nF$. En una conexión pasiva, la resistencia equivalente
    vista por el cable es $R_{\mathrm{th}} = R_S \parallel R_L = 5\kohm$, estableciendo una frecuencia de corte superior de:
    \begin{equation}
      f_c = \frac{1}{2\pi \cdot R_{\mathrm{th}} \cdot C_{\mathrm{cable}}} = \frac{1}{2\pi \cdot 5000\ohm \cdot 3\nF} \approx 10.6\kHz
    \end{equation}
    Esto provocaría atenuación en las frecuencias altas del rango de interés ($10\kHz$ a $20\kHz$).
\end{enumerate}

Para superar estas limitaciones, se requiere intercalar un circuito activo que presente alta impedancia de entrada
para no cargar al micrófono ($Z_i \ge R_S$) y una impedancia de salida reducida ($Z_o \ll R_L$) capaz de excitar la
capacidad de la línea sin atenuación de agudos.

\section{Comparación de Topologías y Elección Justificada}

En el Cuadro~\ref{tab:comparacion_topologias} se resumen las propiedades características de las tres configuraciones
básicas con BJT.

\begin{table}[!ht]
  \centering
  \resizebox{\textwidth}{!}{
  \begin{tabular}{|l|c|c|c|c|c|l|}
    \hline
    \textbf{Topología} & \textbf{Impedancia $Z_i$} & \textbf{Impedancia $Z_o$} & \textbf{Ganancia $A_v$} & \textbf{Ganancia $A_i$} & \textbf{Fase} & \textbf{Aplicación Típica} \\
    \hline
    Emisor Común & Media-alta & Media & Media-alta & Media & $180^\circ$ & Amplificación de uso general \\
    Base Común & Muy baja & Alta & Media-alta & $\approx 1$ & $0^\circ$ & Alta frecuencia / Transimpedancia \\
    Colector Común & Alta & Baja & $\approx 1$ & Media-alta & $0^\circ$ & Separación de etapas (buffer) \\
    \hline
  \end{tabular}
  }
  \caption{Comparativa técnica de las configuraciones con transistor bipolar de juntura.}
  \label{tab:comparacion_topologias}
\end{table}

Analizando cada topología frente al problema:
\begin{itemize}
  \item Configuración Base Común: resulta totalmente inadecuada debido a su impedancia de entrada extremadamente baja
    ($Z_i \approx r_e \approx 25\ohm$). Al conectarla al micrófono ($R_S = 10\kohm$), la señal quedaría casi totalmente
    atenuada en la entrada por divisor resistivo ($25 / (10000 + 25) \approx 0.0025$).
  \item Configuración Emisor Común: provee ganancia de tensión elevada e invierte la fase $180^\circ$. Si bien su
    impedancia de entrada es intermedia, resulta del orden de $1\kohm$ a $2\kohm$, lo cual provocaría una fuerte carga
    sobre los $10\kohm$ del micrófono. Asimismo, la alta ganancia de tensión saturaría el equipo de grabación.
  \item Configuración Colector Común (Seguidor Emisor): ofrece una impedancia de entrada elevada ($Z_i \ge 25\kohm$),
    impidiendo la caída de tensión sobre la resistencia interna del micrófono; presenta muy baja impedancia de salida
    ($Z_o \le 100\ohm$), permitiendo atacar con solvencia la carga de $10\kohm$ y la capacidad parásita del cable de
    $30\unit{\meter}$; y provee una ganancia de tensión unitaria ($A_v \approx 1$) sin inversión de fase ($0^\circ$).
\end{itemize}

Por consiguiente, la configuración Colector Común es la única topología que satisface plenamente las exigencias del
problema de acondicionamiento de señal.
